\section{Introduction}
\label{sec:intro}

Representations carry features that the task does not need. When one of those
features is correlated with a group label, the usual consequence is a biased
prediction, and the usual remedy is to audit or retrain the model that makes it.
That framing assumes the learner is alone. Increasingly it is not: systems are
assembled from many learned agents that read one another's outputs, form some
judgement of how far to rely on each other, and revise accordingly. In such a
population a learner adapts not only \emph{what it predicts} but \emph{whom it
learns from}, and a representation error acquires a second route of influence.
It can reshape the information flow of the whole population.

The collective consequence of that second route is what we study. Online
Bayesian and entropic accounts describe how a single learner should update from
a noisy source
\citep{opper1996online,opper1998bayesian,solla1999optimal,caticha2020entropic};
opinion-dynamics and signed-network models describe how populations polarise and
how sign structures organise
\citep{deffuant2000mixing,castellano2009statistical,heider1946attitudes,cartwright1956structural};
fairness research characterises task-irrelevant and group-correlated features in
representations
\citep{dwork2012fairness,hardt2016equality,barocas2023fairness}. What is missing
is the combination: what a group-correlated representation error does to a
population whose \emph{trust} co-adapts alongside its opinions.

\paragraph{The model, in plain language.}
A population of agents discusses a fixed agenda of issues. Each agent is a
perceptron over a shared embedding of issues, so two agents can agree on some
and disagree on others. Each agent additionally maintains, for every other
agent, an estimate of how unreliable that agent is as a source --- a trust, which
is inferred rather than posited. Both sectors co-adapt through a single update
derived from a Bayesian step on the agent's belief. Every agent also carries an
inert binary label. A fraction $\fd$ of agents extend their representation of an
incoming \emph{message} with one extra coordinate that carries no information
about the issue and depends only on the sender's label, shifting their perceived
agreement by $\pm d$ according to whether the sender shares it.

\paragraph{The mechanism, in one paragraph.}
That shift translates the local opinion--trust flow rigidly, and the translation
moves a separatrix that we can locate exactly. One consonant basin --- ``agree
and trust'' --- grows at the expense of the other, and it grows towards the
in-group and shrinks towards the out-group. Crucially the shift enters the two
sectors asymmetrically: it can flip the \emph{sign} of a trust update, but only
rescales the \emph{magnitude} of an opinion update. Class structure therefore
appears in trust first and reaches opinion only through it. That asymmetry
predicts both the sign and the ordering of the effects we then measure.

\paragraph{Contributions.}
\begin{itemize}\itemsep1pt
  \item \textbf{A formal model of class-correlated representation error}
    (\S\ref{sec:mechanism}). We derive the field \eqref{eq:disc-field} from an
    explicit augmented representation of an interaction, state the assumptions
    under which the reduction is exact, and make precise the sense in which the
    added feature is irrelevant to the task.
  \item \textbf{An exact local analysis} (\S\ref{sec:flow}). The field
    translates the phase portrait rigidly; the separatrix is a line we give in
    closed form; the unbiased flow has a $\mathbb{Z}_2$ symmetry that the field
    breaks. No expansion in $d$ is used anywhere.
  \item \textbf{A parameter-free prediction, and a measurement of what exceeds
    it} (\S\ref{sec:prediction}). A basin-capture argument gives
    $\CcT=\fd\,\Gresp(d)$ with $\Gresp$ closed-form at the model's own
    initialisation. The measured correlation exceeds it by \num{residual}, and
    that residual is the population's own amplification.
  \item \textbf{A replicated map of the collective regimes}
    (\S\ref{sec:regimes}), with regimes defined by a stated classification rule
    against null-calibrated thresholds, boundary uncertainty shown rather than
    asserted, and controls establishing that the effect requires class
    \emph{correlation} and not merely a perturbation of the same size.
  \item \textbf{A temporal result} (\S\ref{sec:agenda}): the agenda complexity
    $\alpha=P/K$ controls whether affective or ideological structure forms first,
    crossing over at $\alpha=\num{crossover}$.
\end{itemize}

\paragraph{A negative result we state plainly.}
The sign change of the trust--class correlation at $d=0$ looks like a transition
and is not one. We show that the affective sector has no linear restoring force
at the class-symmetric state, so no critical field strength for $\CcT$ exists at
any $N$; the sign change is a forced zero-crossing of an odd response. This
explains why an earlier attempt to obtain a critical curve for this model
returned a largest eigenvalue that was identically zero, and it redirects the
search for a genuine threshold to the onset of \emph{ideological} alignment,
where we look for one.

\paragraph{Scope.}
This is a mechanistic model of interacting learners, not a model of people. Its
agents are perceptrons, its groups are arbitrary inert labels, and its
``tolerance'' is a scalar. We draw no conclusions about human societies. What
the results support is a claim about learning systems: a locally well-motivated
update, deployed in a population, can convert an informationally empty label
into a stable architecture of group distrust --- a failure of the interaction
rather than of the training data, and one invisible to any evaluation that
inspects agents one at a time.
